\documentclass[10pt,a4paper]{amsart}
\usepackage[margin=19mm]{geometry}
\usepackage{amsmath,amssymb,mathtools}
\usepackage[scr=rsfs,cal=euler]{mathalfa}
\usepackage{xfrac}
\usepackage[unicode,hidelinks,bookmarks=false]{hyperref}
\newcommand{\ie}{i.e.\ }
\newcommand{\cf}{cf.\ }
\newcommand{\resp}{respectively\ }
\newcommand{\Subsection}{\subsection}
\providecommand{\nobreakdash}{\nobreak\hbox{-}}
\setlength{\emergencystretch}{2em}
\multlinegap0pt
\usepackage{bm}
\usepackage{bbm}
\usepackage{dsfont}
\usepackage{xspace}
\usepackage{cancel}
\usepackage{mathtools}
\usepackage{amsthm}
\makeatletter
\@ifundefined{algorithm}{\newtheorem{algorithm}{Algorithm}}{}
\@ifundefined{theorem}{\newtheorem{theorem}{Theorem}}{}
\renewenvironment{proof}[1]{\par\medskip\noindent\textbf{#1}\enspace\ignorespaces}{\par\medskip}
\makeatother
\title[Resource Allocation Based on Past Incident Patterns]{Resource Allocation Based on Past Incident Patterns}
\author{M.N.M. van Lieshout}
\date{Source: arXiv:2601.16702v1 (CC BY 4.0)}
\begin{document}
\maketitle

\noindent\emph{Source and licence.} Resource Allocation Based on Past Incident Patterns. Version arXiv:2601.16702v1 (CC BY 4.0), distributed under the Creative Commons Attribution 4.0 International licence, as recorded in the arXiv licence metadata for this exact version and shown on the abstract page https://arxiv.org/abs/2601.16702v1. Full rights notice: licenses/arxiv-2601.16702-CC-BY-4.0.txt.\par\smallskip

\noindent\textit{Editorial scope.} Counted unit: the Theorem whose source label is \texttt{t:personnel} in the article (source0.tex lines 537--544, source label \texttt{t:personnel}), delivered together with the article's own complete original proof of it (source0.tex lines 546--669, one \texttt{proof} environment of 124 lines). No mathematics is written, repaired or re-derived: statement and proof are the article's own text line for line. Required context retained uncounted: e:vehicles (lines 300--303); A:pumper (lines 316--335); t:cars (lines 360--367); e:humans (lines 466--470); A:men (lines 488--508). Every definition, display and label of those lines is retained unchanged. Omitted from this package and not redistributed: 1--299, 304--315, 336--359, 368--465, 471--487, 509--536, 545--545, 670--1033. Nothing inside a retained excerpt is omitted or altered: every retained body is delivered exactly as the article's lines are written, so no omission receipt of any kind accompanies this package. Citations: the retained excerpts carry no citation command at all, so no bibliography entry of the article is transcribed and no reference is redistributed. Reference adaptation: none was needed and none was made; every reference inside the delivered unit resolves inside it, so no unresolved reference or citation is produced. Printed slips kept literal: no printed word, symbol, hypothesis or number of the counted statement or of its proof was altered. Layout and class: the article's own class is \texttt{article} with packages including epsf, amssymb, graphicx, natbib; the wrapper is the lane's pinned \texttt{amsart} editorial wrapper at 10pt on A4, which restates the theorem-like environments the retained text needs and adds the article's own macro definitions that the retained text needs (none), with extra wrapper packages (bm, bbm, dsfont, xspace, cancel, mathtools). The delivered numbering is the unit's own. Assets: no retained span contains a figure environment, a graphics-inclusion command, a PDF-page inclusion, a table or a TikZ picture; no image file is copied into this package, the package declares no asset dependency and no image file is redistributed with it. Licence: version arXiv:2601.16702v1 of this article is distributed under the Creative Commons Attribution 4.0 International licence, as recorded in the arXiv licence metadata for this exact version and shown on the abstract page https://arxiv.org/abs/2601.16702v1. A full rights notice is delivered at licenses/arxiv-2601.16702-CC-BY-4.0.txt. Characters: every retained excerpt is pure ASCII, so the delivered engine, XeLaTeX with its default Latin Modern text font, reports no missing character.
\begin{equation}
\min\max_{s \in \mathcal{S}} \frac{\Lambda(s)}{n(s)}
\label{e:vehicles}
\end{equation}over all $n: \mathcal{S} \to \mathbb{N}$ such that 

\begin{algorithm}
Initialise by allocating one resource unit to every station,
i.e. $n_{|\mathcal{S}|}(s) = 1$ for all $s\in \mathcal{S}$.
At iteration $k$, after $k$ resources have already been 
allocated according to $(n_k(s))_{s\in\mathcal{S}}$, 
add one to a station $s^*$ for which 
\[
s^* = \arg\max_{s\in \mathcal{S}} \frac{\Lambda(s)}{n_k(s)}
\]
to obtain 
\[
\left\{ 
\begin{array}{llll}
    n_{k+1}(s) & = & n_k(s) + 1, & \mbox{ if } s = s^* \\
    n_{k+1}(s) & = & n_k(s), & \mbox{ if } s\neq s^*. \\
\end{array} 
\right.
\]
\label{A:pumper}
\end{algorithm}

\begin{theorem} \label{t:cars}
Suppose that $K \geq |\mathcal{S}|$ and that $\Lambda(s) > 0$
for all $s\in \mathcal{S}$. Then, Algorithm~\ref{A:pumper} with
$k=K$ returns a decision rule $n(s)$, $s\in \mathcal{S}$, that 
minimises the maximal average risk (\ref{e:vehicles}) and that 
satisfies the constraints  $1\leq n(s)$ for all $s$ and 
$\sum_{s\in\mathcal{S}} n(s) = K$.
\end{theorem}

\begin{equation}
\min\max_{s \in \mathcal{S}} \frac{\Lambda(s)}
{ \lfloor f(s)/\alpha \rfloor}
\label{e:humans}
\end{equation}

\begin{algorithm}
Initialise by allocating $\alpha$ resource units to every station,
i.e. $f_{\alpha |\mathcal{S}|}(s) = \alpha$ for all $s\in \mathcal{S}$.
At iteration $k$, after $k$ resources have already been 
allocated according to $(f_k(s))_{s\in\mathcal{S}}$, 
add one to a station $s^*$ for which 
\[
s^* = \arg\max_{ f_k(s) < \alpha n(s)} 
  \frac{\Lambda(s)}{ \lfloor f_k(s)/ \alpha \rfloor}
\]
to obtain 
\[
\left\{ \begin{array}{llll}
    f_{k+1}(s) & = & f_k(s) + 1,  & \mbox{ if } s = s^*, \\
    f_{k+1}(s) & = & f_k(s), & \mbox{ if } s\neq s^*.
\end{array} \right.
\]
In case of ties, preference is given to those stations
$s$ for which $f_k(s) \, \rm{mod} \, \alpha$ is largest.
\label{A:men}
\end{algorithm}

\begin{theorem} \label{t:personnel}
Suppose that $n(s) \geq 1$ and $\Lambda(s) > 0$ for all $s\in \mathcal{S}$.
Let $\alpha |\mathcal{S}| \leq K < \alpha \sum_{s\in \mathcal{S}} n(s)$.
Then Algorithm~\ref{A:men} with $k=K$ returns a decision rule $f(s)$,
$s\in \mathcal{S}$ that minimises the maximal average risk (\ref{e:humans}) 
and satisfies the constraints that $\alpha \leq f(s) \leq \alpha n(s)$ 
for all $s\in \mathcal{S}$ and that $\sum_{s\in \mathcal{S}} f(s) = K$.
\end{theorem}

\begin{proof}{}
The greedy algorithm is initialised with 
$f_{\alpha |\mathcal{S}}(s) = \alpha$ and 
continues until $K$ resources have been allocated. 
Additionally, in each step, the addition of one resource
is constrained to $f_k(s) < \alpha  n(s)$.
Hence  the constraints are satisfied. 

The optimality proof proceeds by contradiction.  If $(f_k(s))_{s\in\mathcal{S}}$
is not optimal, there exists an optimal solution $(f(s))_{s\in\mathcal{S}}$ that is strictly better, i.e.,
\[
    \max_{s\in\mathcal{S}} \frac{\Lambda(s)}
    { \lfloor f(s) / \alpha \rfloor }
    < \max_{s\in\mathcal{S}} \frac{\Lambda(s)}
     { \lfloor f_k(s) / \alpha \rfloor }.
\]

By the tie-breaker rule, the 
output of Algorithm~\ref{A:men}, $f_k$, contains at 
most one station with slack, which has size
$K - \alpha\lfloor K/\alpha \rfloor \leq \alpha -1$,
and for all other stations, $f_k$ is a multiple of $\alpha$. 

For $f$, if the combined slack is $\alpha$ or more, pick any 
$s_0 = \arg\max_s \Lambda(s) /$ $ \lfloor 
( f(s) / \alpha ) \rfloor$ with 
$\alpha \lfloor f(s_0) / \alpha \rfloor + \alpha \leq n(s_0)$.
If the maximiser $s_0$ is unique, 
allocating units out of the slack to station $s_0$ to 
fill up an additional vehicle 
would lead to a smaller maximal risk per crew, thus 
contradicting the assumption that $f$ is optimal. In case
of a tie, this action would leave the maximal risk per crew 
unchanged and we repeat this procedure
until the total slack is smaller than 
$\alpha$ or no maximiser $s_0$ can accommodate an extra crew.
In the latter case, filling up  remaining stations to
full capacity  as long as there at least $\alpha$ units
left, does not change the maximal risk per vehicle and
the procedure terminates because of the scarcity assumption
$K < \alpha  \sum_{s} n(s)$. 

In conclusion, we may assume without loss 
of generality that also in $f$ the total slack is 
$K - \alpha \lfloor K/\alpha \rfloor \leq \alpha - 1$. 
Set $\tilde K = \alpha \lfloor K/ \alpha \rfloor$. Then 
removing the slack in both $f_k$ and $f$, the remnants 
$\tilde f_k$ and $\tilde f$ now satisfy 
$\sum_s \tilde f(s) = \sum_s \tilde f_k(s) = \tilde K$.
The risks per crew remain unchanged, so $\tilde f$ 
is optimal. Moreover, all $\tilde f(s)$ and $\tilde f_k(s)$ are 
multiples of $\alpha$. Without loss of generality, we may 
now choose $\tilde f$ such that 
\[
\sum_{s\in \mathcal{S}} \left|
 \frac{\tilde f(s)}{\alpha} 
 - \frac{\tilde f_k(s)}{\alpha} 
\right|
\]
is smallest. 

The proof now proceeds along the same line as that of 
Theorem~\ref{t:cars}. If $\tilde f$ and $\tilde f_k$ are not 
the same, there exist $s$ and $s^\prime$ in $\mathcal{S}$
such that $\tilde f_k(s) > \tilde f(s)$ and 
$\tilde f_k(s^\prime) < \tilde f(s^\prime)$. 
We can therefore create a new solution $\hat f$ from 
$\tilde f$ by adding $\alpha$ resources to $s$ and 
withdrawing $\alpha$ units from $s^\prime$. Note that 
$\hat f(\cdot)$ satisfies the constraints and has
smaller $\ell_1$-distance to $\tilde f_k(\cdot)$.  

Observe that
\[
\frac{\Lambda(s^\prime)}
    { ( \tilde f(s^\prime)-\alpha) / \alpha }
    \leq \frac{\Lambda(s^\prime)}
     { \tilde f_k(s^\prime) / \alpha }
\mbox{ and }
   \frac{\Lambda(s)}
   { (\tilde f_k(s) -\alpha)/ \alpha}
   \leq  \frac{\Lambda(s)}
    {\tilde f(s) / \alpha }.
\]
Since the greedy algorithm has selected $\tilde f_k(s)$ but not
$\tilde f_k(s^\prime) + \alpha$,
\[
\frac{\Lambda(s^\prime)}
    { \tilde f_k(s^\prime) / \alpha  }
     \leq  \frac{\Lambda(s)}
    { (\tilde f_k(s) - \alpha)/ \alpha  }
\]
and therefore
\[
\frac{\Lambda(s^\prime)}
    {  ( \tilde f(s^\prime)-\alpha) / \alpha  }
    \leq \frac{\Lambda(s)}
    {  \tilde f(s) / \alpha  }.
\]
Hence 
\[
\frac{\Lambda(s^\prime)}
    { \lfloor \hat f(s^\prime) / \alpha \rfloor } =
\frac{\Lambda(s^\prime)}
    {  ( \tilde f(s^\prime)-\alpha) / \alpha } \leq 
    \frac{\Lambda(s)}
    {  \tilde f(s) / \alpha }
    \leq \max_{t\in \mathcal{S}} \frac{\Lambda(t)}
    { \lfloor \tilde f(t) / \alpha \rfloor }
\]
and
\[
\frac{\Lambda(s)}
    { \lfloor \hat f(s) / \alpha \rfloor } =
\frac{\Lambda(s)}
    {  ( \tilde f(s)+ \alpha) / \alpha} \leq 
    \frac{\Lambda(s)}
    {  \tilde f(s) / \alpha }
    \leq \max_{t\in \mathcal{S}} \frac{\Lambda(t)}
    { \lfloor \tilde f(t) / \alpha \rfloor },
\]
leading to a contradiction with the assumption that $\tilde f$
is an optimal allocation rule closest to $\tilde f_k$. The proof is complete.
\end{proof}

\end{document}
